\par\Needspace{6\baselineskip}\medskip
\hypertarget{native-155-article111}{}\pdfbookmark[2]{111.}{article-111}
\begin{center}111.\end{center}
We next suppose that the two positions are given by their distances from three planes intersecting at right angles at the Sun. Let these distances be denoted by $x$, $y$, and $z$ for the first position and by $x'$, $y'$, and $z'$ for the second. Suppose that the third plane is the ecliptic itself and that the positive poles of the first and second planes lie at longitudes $N$ and $90^0+N$. Then, by Article 53, if the two radius vectors are denoted
\SourceLeaf{156}by $r$ and $r'$, we have
\[
\begin{aligned}
x&=r\cos u\cos(N-\AscendingNode)+r\sin u\sin(N-\AscendingNode)\cos i\\
y&=r\sin u\cos(N-\AscendingNode)\cos i-r\cos u\sin(N-\AscendingNode)\\
z&=r\sin u\sin i\\
x'&=r'\cos u'\cos(N-\AscendingNode)+r'\sin u'\sin(N-\AscendingNode)\cos i\\
y'&=r'\sin u'\cos(N-\AscendingNode)\cos i-r'\cos u'\sin(N-\AscendingNode)\\
z'&=r'\sin u'\sin i.
\end{aligned}
\]
It follows that
\[
\begin{aligned}
zy'-yz'&=rr'\sin(u'-u)\sin(N-\AscendingNode)\sin i\\
xz'-zx'&=rr'\sin(u'-u)\cos(N-\AscendingNode)\sin i\\
xy'-yx'&=rr'\sin(u'-u)\cos i.
\end{aligned}
\]
Combining the first and second formulas gives $N-\AscendingNode$ and $rr'\sin(u'-u)\sin i$. From these and the third formula we obtain $i$ and $rr'\sin(u'-u)$.

If the position corresponding to the coordinates $x'$, $y'$, and $z'$ is taken to be the later one, $u'$ must be greater than $u$. If it is also known whether the angle swept out about the Sun between the first and second positions is less or greater than two right angles, $rr'\sin(u'-u)\sin i$ and $rr'\sin(u'-u)$ must be positive in the first case and negative in the second. Then $N-\AscendingNode$ is determined unambiguously, and the sign of $xy'-yx'$ decides whether the motion is direct or retrograde. Conversely, if the direction of motion is known, the sign of $xy'-yx'$ determines whether $u'-u$ is to be taken as less or greater than $180^0$. If both the direction of motion and the character of the angle swept out about the Sun are wholly unknown, it is clear that the ascending node cannot be distinguished from the descending node.

It is readily seen that, just as $\cos i$ is the cosine of the inclination of the orbital plane to the third plane, so $\sin(N-\AscendingNode)\sin i$ and $\cos(N-\AscendingNode)\sin i$ are the cosines of its inclinations to the first and second planes, respectively. Also, $rr'\sin(u'-u)$ expresses twice the area of the triangle enclosed by the two radius vectors, and $zy'-yz'$, $xz'-zx'$, and $xy'-yx'$ express twice the areas of the projections of the same triangle onto the respective planes.

Finally, it is clear that any other plane can serve as the third plane in place of the ecliptic, provided that all quantities defined by their relations to the ecliptic are referred in the same way to whichever third plane is chosen.
